\documentclass[aps,prd,twocolumn,superscriptaddress,nofootinbib]{revtex4-2}
\usepackage{amsmath,amssymb}
\usepackage{graphicx}
\usepackage{hyperref}
\usepackage{booktabs}
\usepackage{orcidlink}

\begin{document}

\title{One Register:\\ A Common Origin for the Granularity of Hilbert Space and of Space}

\author{Andrew Korytko\,\orcidlink{0009-0005-4569-2228}}
\email{andrew@alphalatitude.com}
\affiliation{AlphaLatitude Inc., P.O. Box 64341, Sunnyvale, California 94088, USA}

\date{\today}

\begin{abstract}
In the paper on the origin of gravitation, we postulated separately that Hilbert space and space are granular with one universal granularity $L$, extending Palmer's Rational Quantum Mechanics. We derived Newton's constant and Einstein's equations, and pinned $L$ at two scales, galactic and cosmological. Here the two axioms become one postulate: a qubit is a record of marks, one at each cell of a ring of $L$ Planck cells, and a state of $N$ qubits is $N$ records on the same cells, read jointly. Its fraction $m/L$ of $+1$ marks is a Born probability and its cyclic position a phase, one step of phase being one cell, a full turn one circuit. The ring is the register of the title. Placing the ring on a great circle of the cosmological horizon, as before, gives $L = 2\pi R_\Lambda/\ell_P \approx 6\times10^{61}$. Momentum is quantized at $\hbar/R_\Lambda$; a state spread over all $2^N$ outcomes of $N$ qubits needs cells for each, $2^N \le L$, at most $\lfloor\log_2 L\rfloor = 205$, whatever realizes them. Along a direction, light moves one cell per tick, velocity is $v/c = 2m/L - 1$, and matter at rest reverses every quarter Compton period, its clock Palmer's hidden global phase; several particles in several planes need three dimensions, a count derived elsewhere. Predictions: the qubit ceiling, a random circuit and its inverse failing beyond 205 logical qubits, constant $\Lambda$, no energy-dependent light speed, and, unlike collapse models, no intrinsic matter-wave decoherence.

\end{abstract}

\keywords{Rational quantum mechanics; discrete Hilbert space; finite quantum mechanics; de Sitter horizon; Bell's theorem}

\maketitle

\section{Introduction}
\label{sec:intro}

A qubit is a point on the Bloch sphere,
\begin{equation}
|\psi\rangle = \cos\tfrac{\theta}{2}\,|1\rangle + e^{i\phi}\sin\tfrac{\theta}{2}\,|{-1}\rangle ,
\label{eq:qubit}
\end{equation}
and in standard quantum mechanics every point is a state. Rational Quantum Mechanics (RaQM)~\cite{palmer2020,palmer2026,palmer2026b,palmer2026c} keeps the Schr\"odinger equation and the Born rule but admits only those points, and only those measurement bases, for which
\begin{equation}
\cos^2\tfrac{\theta}{2} = \frac{m}{L},\qquad \frac{\phi}{2\pi} = \frac{n}{L},
\label{eq:raqm}
\end{equation}
with $m\in\{0,\dots,L\}$, $n\in\{0,\dots,L-1\}$, and $L$ a large integer. A quantum state is then not a point on a continuum but a point of a very fine grid, and $L$ is how fine the grid is. Palmer represents such a state as a string of $L$ bits: the fraction of $+1$ bits is $\cos^2(\theta/2)$, and rotating the string by $n$ places is the phase~\cite{palmer2026}. Niven's theorem~\cite{niven1956} says that $\cos\alpha$ and $\alpha/2\pi$ are simultaneously rational only for a handful of angles, so Eq.~(\ref{eq:raqm}) forbids the pair of counterfactual settings a Bell inequality needs, which is Palmer's Impossible Triangle Corollary, and Bell's theorem no longer implies nonlocality~\cite{palmer2020,palmer2026b}.

A state of $N$ qubits spread across its Hilbert space has $2^N$ outcomes, and with probabilities read as frequencies over $L$ trials there is a largest $N$ for which such a state can be written at all, $N_{\max}\approx\log_2L$~\cite{palmer2026}; the bit count $2^{N+1}-2 \le NL$, Eq.~(10) of Ref.~\cite{palmer2026c}, is a weaker condition. In Palmer's version $L$ is a property of each qubit, set by its mass and energy, and ranges from $10^{64}$ to $10^{109}$.

In a previous paper~\cite{korytko2026} we made two changes and called them axioms: $L$ is a single constant of nature, and one step of quantum phase, $2\pi/L$, is one Planck length $\ell_P = \sqrt{\hbar G/c^3}$ on a great circle of the cosmological horizon at $R_\Lambda = \sqrt{3/\Lambda} = 1.66\times10^{26}$~m~\cite{planck2018},
\begin{equation}
\ell_P = \frac{2\pi R_\Lambda}{L} .
\label{eq:ident}
\end{equation}
From the two, Newton's constant follows by Verlinde's argument~\cite{verlinde2011} with the pixel fixed, Einstein's equations by Jacobson's~\cite{jacobson1995}, rotation curves flatten at $a_0 = c^2/L\ell_P$, and $G\Lambda = 12\pi^2c^3/\hbar L^2$. Milgrom's $a_0$~\cite{milgrom1983,mcgaugh2016} gives $L = 4.6\times10^{61}$ and the observed $\Lambda$ gives $6.4\times10^{61}$, agreeing to a factor 1.4; both lie below every per-qubit value in Palmer's range.

Equation~(\ref{eq:ident}) equates the granularity of two things introduced as separate, Hilbert space and space. Ref.~\cite{korytko2026} deferred to a separate paper the unification of its two axioms and the extension of the first to continuous variables; this is that paper. It takes the identification, with the horizon withheld, as its only postulate and draws out what it contains: a step that turns a phase and moves a length at once, whose repetitions close on a ring of $L$ Planck cells. The first axiom follows, the rational grid with $L$ universal; the second, placing the ring on the cosmological horizon, is taken over here, Sec.~\ref{sec:horizon}, where it gives $L$ its value, and Ref.~\cite{korytko2026d} derives it from the same discreteness with one postulate more, that a direction is the state of one qubit. Palmer's phase step is what the step does to the slowest wave the closure admits, and the count of steps that close is $L$. Read for a qubit, the rotation is a phase; read for a particle, a translation; for light the two are one thing at every Planck tick.

Section~\ref{sec:register} states the postulate and derives from it the first axiom of Ref.~\cite{korytko2026}. Sections~\ref{sec:algebra} to \ref{sec:horizon} draw what follows from the step alone, with no direction chosen: its algebra and momenta, the capacity of the register, and, with the ring on the horizon, the quantities the horizon fixes. Section~\ref{sec:direction} fixes a direction, where light, matter, and velocity appear and position in whole cells is one-dimensional; Sec.~\ref{sec:beyond} says how far that reaches, several particles in several planes needing a three-dimensional space that the paper does not construct and whose dimension Ref.~\cite{korytko2026d} derives. Section~\ref{sec:pred} states what is predicted and what would falsify it, and Sec.~\ref{sec:cost} what the placement assumes and what no register holds.

\section{One Discreteness}
\label{sec:register}

\textbf{Postulate.} \emph{The granularity of state space and the granularity of space are one discreteness. Every qubit is a record of marks, one to each cell of a ring of $L$ cells, each cell one Planck length and one Planck time, each mark $+1$ or $-1$: the fraction of $+1$ marks is a Born probability and the cyclic position of the record is its phase, so that one step of quantum phase is one cell and a full turn of phase is one circuit of the ring. The global phase of a state is hidden: no law depends on it, and the dynamics does not expose it.}

Phase is cyclic and space is not, so it is the phase that closes. A full turn returns a phase to itself, and the postulate says what that turn costs in cells. $L$ is their number: the cells in the ring, the marks in a record, and the steps in a full turn, since a record of $L$ marks has $L$ cyclic positions and one step carries it to the next. The second sentence read for the marks is the rationality condition of the first axiom of Ref.~\cite{korytko2026}, Sec.~\ref{sec:grid}, with the universality clause a consequence, Sec.~\ref{sec:counting}; read for the phase it is that paper's second axiom with the horizon withheld and the Planck time added; the third sentence is new. Nothing else is assumed. Where the ring lies, and so how large $L$ is, the postulate does not say: Sec.~\ref{sec:horizon} places it on a great circle of the cosmological horizon, the second axiom of Ref.~\cite{korytko2026}, which Ref.~\cite{korytko2026d} derives from this postulate and one more.

\textbf{The postulate in plain words.}

\textbf{One discreteness.} Two things in physics are grainy for reasons usually thought unrelated. Quantum phase, on Palmer's account, comes in a smallest step rather than varying continuously; and space, on holographic accounts~\cite{thooft1993,susskind1995}, carries at most one degree of freedom per Planck cell of its boundary. The postulate is that these are not two facts but one, and that the single grain, looked at one way, is a turn of phase and, looked at the other, a movement through space.

\textbf{A qubit is a record of marks on the cells.} What a qubit is, on the postulate, is a pattern of $+1$ and $-1$ marks around a ring of $L$ Planck cells, one mark to a cell. The fraction of $+1$ marks is its Born probability and the position of the pattern around the circle is its phase; a qubit that could not be so written does not exist, and neither does a measurement whose outcomes could not be. A state of $N$ qubits is $N$ records on the same cells, $NL$ marks read jointly, Sec.~\ref{sec:records}, and is not a qubit.

\textbf{One step of phase is one cell.} The Planck length, $1.6\times10^{-35}$~m, is the shortest distance with operational meaning, and the cell is space's grain at that scale. The clause says the smallest turn of phase and the smallest cell of space are the same grain counted two ways, so that asking how many steps make a turn is asking how many cells make the ring. It also says the grain is a Planck time as well as a Planck length, $\ell_P = ct_P$, which is the tick of the dynamics of Sec.~\ref{sec:direction}.

\textbf{A full turn of phase is one circuit of the ring.} Phase returns to itself after $2\pi$; the clause says that return corresponds to going once around the ring. Nothing here says space closes on itself. What closes is the phase, and the ring is what its closing costs. Because the expansion of the universe is accelerating, every observer has a sphere of radius $R_\Lambda\approx1.7\times10^{26}$~m beyond which nothing can ever be seen, and Sec.~\ref{sec:horizon} places the ring on a great circle of it, the largest circle it has.

\textbf{The global phase of a state is hidden.} Every quantum state carries an overall phase that no measurement reads; that much is ordinary quantum mechanics. The postulate adds that it stays hidden: the phase advances, at the rate Sec.~\ref{sec:matter} gives, but no law depends on it and nothing the dynamics does exposes it. A state of several records read jointly, Sec.~\ref{sec:records}, has one global phase; the relative arrangement of its records is its correlations, and it is data. Without this clause the gauge argument of Sec.~\ref{sec:cost} is incomplete, since a step with no privileged starting place still leaves open whether a state's phase is referred to anything outside it, and the clause says it is not.

\subsection{What follows from counting}
\label{sec:counting}

Four things follow before the ring is placed anywhere, and with them the first axiom of Ref.~\cite{korytko2026}: the universality it stipulated; the phase step; the momenta; and, in Sec.~\ref{sec:grid}, its rationality condition, Eq.~(\ref{eq:raqm}).

\emph{Universality.} $L$ is the same for every system. The cell is a Planck length, and no system has a cell of its own; one step of phase is one cell, so the phase step is the same for every state; and a full turn is $2\pi$ for every state, so the number of steps in a turn is one number. A system with its own $L$ would need its own cell. This is the universality clause of the first axiom of Ref.~\cite{korytko2026}.

The postulate speaks of \emph{the} granularity of space, so it assumes one cell for everything; universality of $L$ is then a consequence and not a further assumption. A Planck length that differed from system to system is a proposal nobody makes; an $L$ that differs from system to system is exactly Palmer's, set by a system's mass and energy and ranging from $10^{64}$ to $10^{109}$, and is what the first axiom of Ref.~\cite{korytko2026} had to contradict outright. Identifying the two granularities makes that contradiction unnecessary: Palmer's $L$ cannot vary with mass and energy while the cell it is identified with does not.

\emph{The phase step.} $L$ steps make a full turn, so one advances a phase by
\begin{equation}
\Delta\phi = \frac{2\pi}{L},
\label{eq:step}
\end{equation}
which is Palmer's phase step, and his operator $\zeta$, a cyclic permutation of an $L$-bit string, is one application of the step~\cite{palmer2026}, written $X$ in Sec.~\ref{sec:algebra}.

\emph{Momenta.} A plane wave $e^{ipx/\hbar}$ advances in phase by $p\ell_P/\hbar$ across one cell. The phase difference between two cells is a whole number of steps, the postulate's grid, so by Eq.~(\ref{eq:step}) that advance is $2\pi\kappa/L$ for an integer $\kappa$, and
\begin{equation}
p_\kappa = \frac{2\pi\hbar\kappa}{L\ell_P} = \frac{h\kappa}{L\ell_P},
\label{eq:pquant}
\end{equation}
and momentum is quantized in units of $h/L\ell_P$; what is cut off is wavelength and not space, for no wave is longer than the ring. The wave with $\kappa = 1$ advances by exactly one step per cell: Palmer's unit is the step of the longest wavelength there is, and its wavelength is the ring, $L\ell_P$.

\subsection{The grid}
\label{sec:grid}

A record has $L$ marks, so its count of $+1$ marks runs over $L+1$ values, and read for a qubit that count is the Born probability. A qubit's latitude cuts the Bloch sphere into two caps whose fractional areas are $\cos^2(\theta/2)$ and $\sin^2(\theta/2)$; partition the sphere into $L$ equal-area bands, whose $L+1$ boundaries are the values of the count, and a state $m$ bands from the $|{-1}\rangle$ pole has
\begin{equation}
\cos^2\tfrac{\theta}{2} = \frac{m}{L}, \qquad \cos\theta = \frac{2m}{L} - 1 \in \mathbb{Q},
\label{eq:count}
\end{equation}
since equal areas are equal steps in $\cos\theta$ by Archimedes' theorem. The shift is the azimuth: $L$ steps make a full turn, so
\begin{equation}
\frac{\phi}{2\pi} = \frac{n}{L}.
\label{eq:shift}
\end{equation}
Together these are Eq.~(\ref{eq:raqm}), Palmer's rationality condition, with $L$ the cell count of the ring, universal by Sec.~\ref{sec:counting}: it follows from the postulate rather than standing beside it.

The two conditions are independent, $\theta$ and $\phi$ being separate coordinates, so a permitted count may be combined with any permitted shift and the permitted set has of order $L^2$ points. Niven's theorem~\cite{niven1956,jahnel2010} does not bear on that pairing. It constrains a single angle, forbidding $\cos\alpha$ and $\alpha/2\pi$ to be rational together at more than eight angles, and in RaQM what it constrains is the angle between two measurement settings.

Equation~(\ref{eq:raqm}) is read on the Bloch sphere in two ways: as the positions a state may occupy, and as the directions a measurement may be set to. Palmer's Bell results are of the second kind, and turn on a setting being freely choosable; they are inherited here and not re-derived. The register's own qubit is of the first kind, and Sec.~\ref{sec:direction} identifies it.

\subsection{Records}
\label{sec:records}

A qubit is a \emph{record}: $L$ marks, the count $m$ of $+1$ marks giving Eq.~(\ref{eq:count}) and the cyclic position giving Eq.~(\ref{eq:shift}), and the $L$ cells it is written on are the ring's, one mark to a cell. The cells are the orbit of the step: applying the step, written $X$ in Sec.~\ref{sec:algebra}, carries each to the next and $L$ applications return it, so they are the $L$ elements of $\mathbb{Z}_L$. The $L$ cells with the records written over them are the \emph{register}.

Several records are several strings over the same $L$ cells, as several functions share one domain without sharing values, so a cell carries one mark of each and nothing is used up by the writing. The cells are a domain and not a medium, which is why the number of records is not bounded by $L$. The construction takes the records to be written on one ring; how records on different rings, the directions of Ref.~\cite{korytko2026d}, are read jointly is the bulk state of Sec.~\ref{sec:beyond}, which this paper does not construct.

Correlations are written the way Palmer writes them~\cite{palmer2026c}, by reading records jointly, cell by cell. Read the first record plainly and its fraction of marks is its marginal. Read the second on the cells where the first reads $+1$ and on those where it reads $-1$: each set is a \emph{block}, and its length is $L$ times the probability of the branch it reads. A block's fraction of marks is the conditional, on the grid of the block's length. The position of its pattern on the ring, taken against the same origin as every other pattern, is the conditional phase: the phase of the record on that branch, which quantum mechanics reads by selecting the branch and interfering the rest. Two records have four conditional phases with one relation among them, and each order reads three, which is what the state has; for $N$ records each order reads $2^N - 1$. Read the $k$th on the cells reading each branch of the records before it. The joint probability of an outcome is then the fraction of cells reading it, a multiple of $1/L$, and it and the conditional phases are the same in every order, so the order in which the records are read is a choice of coordinates. Equation~(\ref{eq:raqm}) holds of each record as a whole, its total count the marginal. The arrangement of a record read jointly is therefore not the arrangement of the same record read alone, and Palmer's single-run form is the latter.

A block need not be a branch. It is a set of branches on whose cells the same fraction is written, to the resolution each branch's cells allow. The branches it gathers are those on which the earlier records satisfy one value of whatever the conditional depends on. For a product state that is nothing, and the $k$th record is one block. For a stabilizer state it is a parity of earlier marks. For a state whose correlations decay it is the marks within a correlation length, since a record whose effect on a conditional is below one part in the block's length does not segment it on the grid. For a state whose every conditional depends on every earlier record with no rule the marks can be read by, it is the branch itself, one block for each. Two branches whose conditionals are merely close are not one block: a conditional that depends on the whole branch is the count on that branch's own cells, and pooling cannot supply it. Every joint distribution is written this way, entangled ones no less than product ones, so $N$ records carry an arbitrary joint distribution of $N$ qubits and the correlation lies in which cells are read rather than in any structure beyond the marks. The blocks shorten with depth, to $L$ times a branch's probability, of order $L/2^{k-1}$ for a state spread evenly, and that is what Sec.~\ref{sec:capacity} counts. Nothing here needs $L$ to be even or a power of two: when $2^{k-1}$ does not divide $L$, the nearest permitted state has blocks that differ by one cell, and every branch still needs a cell, so the bound of Sec.~\ref{sec:capacity} is unchanged. Which arrangements of marks yield a given set of conditional phases is not constructed here; the counting of Sec.~\ref{sec:capacity} needs only the cells. A relabeling of the cells common to every record leaves every count and every relative position unchanged, so the origin is a choice of coordinates, the hidden phase of the postulate; what is data is the fraction of marks in each block and where the patterns sit relative to one another.

\subsection{What is assumed and what is taken}
\label{sec:provenance}

Assumed are $\hbar$, $c$, the numerical value of $\ell_P$, from which $G = \ell_P^2c^3/\hbar$ is derived in Ref.~\cite{korytko2026} rather than the reverse, and the postulate. Taken from Ref.~\cite{korytko2026} is the placement of the ring on the cosmological horizon, Sec.~\ref{sec:horizon}, which that paper postulated and Ref.~\cite{korytko2026d} derives. Nothing is taken from RaQM: its rationality condition, Eq.~(\ref{eq:raqm}), and its reading of the Born probability as a frequency of marks~\cite{palmer2026} are the postulate's second sentence read for the marks, Sec.~\ref{sec:grid}. Taken from quantum mechanics are the Hilbert space and, for matter, the Dirac mass term, Eq.~(\ref{eq:dirac}). The dynamics of Sec.~\ref{sec:direction} is the postulate's own, since the step is a Planck length and a Planck time and light crosses one per tick.

The integer $L$ plays three roles, and the postulate is why they are one number rather than three: it is the order of the step, so the length of a record; it is the number of cells around the ring and, once the ring is placed on the horizon in Sec.~\ref{sec:horizon}, of pixels around a great circle, equivalently the number of Planck times light takes to go once around; and, once a direction is chosen in Sec.~\ref{sec:direction}, it is the dimension of a free particle's orbital Hilbert space.

\section{The Algebra of the Step}
\label{sec:algebra}

Under the Schr\"odinger equation the phase of a mode of energy $E$ advances by $E\,dt/\hbar$. Palmer's $\zeta$ advances it by $2\pi/L$. Schr\"odinger evolution, in Palmer's representation, is therefore the string of a mode of energy $E$ rotating at
\begin{equation}
\dot n = \frac{E\,L}{2\pi\hbar}\ \text{steps per second} ,
\label{eq:rate}
\end{equation}
and in a superposition each mode's string turns at its own rate, the differences between rates being the phases that show; a single mode's rotation is the hidden phase of Sec.~\ref{sec:matter}.

One cell is one Planck length and one Planck time, $\ell_P = c\,t_P$, by the postulate; the step has one unit and $c$ converts it. Write $X$ for the step. On the $L$ cells of a record it is the shift $|j\rangle\mapsto|j+1\rangle$, and since a cell is a Planck length it is a translation,
\begin{equation}
X|j\rangle = |j+1\rangle = e^{-i\hat p\,\ell_P/\hbar}|j\rangle ,
\label{eq:X}
\end{equation}
which defines momentum with no continuum limit: momentum is what generates the step. Its eigenstates are the plane waves
\begin{equation}
|\tilde\kappa\rangle = \frac{1}{\sqrt L}\sum_j e^{2\pi i j\kappa/L}|j\rangle ,
\label{eq:pk}
\end{equation}
with $-L/2 < \kappa \le L/2$, whose eigenvalues return Eq.~(\ref{eq:pquant}) in operator form: the phase advances by $2\pi\kappa/L$ per cell, a permitted phase with $n = j\kappa \pmod L$, so $p_\kappa = h\kappa/L\ell_P$ as before. The mode $\kappa = 1$, of wavelength $L\ell_P$, the ring, carries the smallest momentum there is. The states whose phase advances by a permitted amount per pixel are the momentum eigenstates; the phase step is the momentum quantum $h/L\ell_P$; position and momentum each take exactly $L$ values; momentum is bounded at the zone edge $\pi\hbar/\ell_P$; and the phase-space cell is $L\ell_P\cdot(2\pi\hbar/\ell_P)/L = h$.

\section{Capacity}
\label{sec:capacity}

Let $Z|j\rangle = e^{2\pi ij/L}|j\rangle$. Then $ZX = e^{2\pi i/L}XZ$, the finite Weyl relation, and Eqs.~(\ref{eq:X})--(\ref{eq:pk}) are Schwinger's finite quantum mechanics on $\mathbb{Z}_L$~\cite{weyl1931,schwinger1960,vourdas2004}. A state on $N_x$ cells occupies at least $L/N_x$ momenta~\cite{donoho1989}, so $\Delta x\,\Delta p \ge (N_x\ell_P)\cdot(L/N_x)(2\pi\hbar/L\ell_P) = h$, and in variance form it is $\hbar/2$ for $\ell_P\ll\sigma\ll L\ell_P$~\cite{massar2008}. Palmer reaches the same bound by averaging the spin relation over points uniform in $\cos\theta$~\cite{palmer2026c}; that measure is the count and this bound comes from the shift, and the two derivations are the register's two operations applied to one relation.

The Hilbert space spanned by the $L$ cells has dimension $L$. The capacity of the register is fixed by room and not by bits. Read jointly, $N$ records give at each cell one of the $2^N$ joint outcomes, and the probability of an outcome is the fraction of cells reading it. A state spanning its Hilbert space gives every outcome a nonzero probability, so every outcome occupies at least one cell, and
\begin{equation}
2^N \le L, \qquad N_{\max} = \lfloor \log_2 L \rfloor ,
\label{eq:nmaxspan}
\end{equation}
205 at $L_{\rm cos} = 6.4\times10^{61}$ and 204 at $L_{\rm gal} = 4.6\times10^{61}$, the cosmological and galactic determinations of Sec.~\ref{sec:horizon}. Branch by branch it is the same condition: a conditional strictly between 0 and 1 needs two cells, and the last record has $2^{N-1}$ branches. The bound is exact, and it is a ceiling any representation must respect and not a construction that attains it, which is the direction a falsifier needs. Its softness is in resolution rather than in the integer: a branch of two cells holds a conditional of one half and one of two phases, so a state whose deepest conditionals and phases must be told apart to one part in $r$ needs $r$ cells to a branch, $2^{N-1}r \le L$, and the ceiling falls by $\log_2(r/2)$. With $L$ known to a factor 1.4 besides, the ceiling sits between 203, for three bits of resolution, $r = 8$, and 205. What the argument fixes is the scale, $\log_2 L$, and that the ceiling is the same for every physical realization; the integer within it is not fixed to better than a unit, and nothing below turns on which.

It is the number of blocks and not the number of records that Eq.~(\ref{eq:nmaxspan}) constrains, and that difference is its entire content. A product state has one block per record, and the bound asks nothing of it at any $N$, as it must not for $|{+}\rangle^{\otimes N}$. A stabilizer state has at most two blocks per record, one for each value of a parity of earlier marks, its conditional 0, 1, or one half, and a state whose correlations decay with a length $\xi$ has of order $2^{\xi}$. So a fault-tolerant code, a Slater determinant, a condensate, a nucleus, or a crystal clears the bound by scores of orders of magnitude, and the register limits neither how much matter exists nor the size of a code. What fills the last record with $2^{N-1}$ blocks is scrambling, a state whose every conditional depends at order one on every earlier record, which is what a random circuit produces once it anticoncentrates, at a depth logarithmic in $N$~\cite{dalzell2022}. No open system holds such a state, since a system entangled with its surroundings is described by a mixed state and a mixture costs no room beyond what its members need; a closed chaotic system isolated for a scrambling time does, which is what a quantum computer is. The one natural object that appears to exceed the bound, a black hole, is taken up in Sec.~\ref{sec:cost}.

Palmer's bit count, $2^{N+1}-2 \le NL$, gives 212 at $L_{\rm cos}$~\cite{palmer2026c}. It charges the $NL$ marks one bit for each real parameter, but what stops a state being written is not a shortage of bits: a record of $L$ cells holds at most $L$ bits whatever is on it, and a record whose $L$ blocks hold one cell each uses every bit and spans nothing, its conditionals all 0 or 1. The constraint is room, two cells to a branch, so the 205 stands and the 212 does not. Ref.~\cite{korytko2026} gave $N_{\max}\approx205$ from $\log_2 L$ and quoted the prediction as about 200; it also called the bit count the exact criterion and inverted a measured $N_{\max}$ into $G$ through it. The inversion keeps its form, Eq.~(\ref{eq:nmaxspan}) bracketing $L$ between $2^{N_{\max}}$ and $2^{N_{\max}+1}$. Equation~(\ref{eq:nmaxspan}) counts outcomes and not the systems that produce them, so a Hilbert space of any dimension $D$ spread across its outcomes obeys $D \le L$, for bosonic modes as for qubits.

\section{The Horizon}
\label{sec:horizon}

The cells close in a circuit of circumference $L\ell_P$ and radius $L\ell_P/2\pi$. The second axiom of Ref.~\cite{korytko2026} places it on a great circle of the cosmological horizon,
\begin{equation}
R_\Lambda = \frac{L\ell_P}{2\pi}, \qquad \Lambda = \frac{3}{R_\Lambda^2} = \frac{12\pi^2}{L^2\ell_P^2},
\label{eq:closure}
\end{equation}
which is Eq.~(\ref{eq:ident}). The postulate does not contain the placement and this paper does not derive it. Ref.~\cite{korytko2026d} does, from the postulate together with one more, that a direction is the state of one qubit, under which the directions of space are a qubit's sphere of states, a two-sphere of this radius. Here the placement is taken over, and it is what gives $L$ its value, $2\pi R_\Lambda/\ell_P$. The identification is one of structure, the same $\mathbb{Z}_L$ under a qubit's phase and a horizon great circle; this paper does not need a laboratory spin's azimuth to be a place on the sky, an identification Ref.~\cite{korytko2026d} makes as its second postulate. Section~\ref{sec:counting} gives the axiom its content: the circumference is the wavelength of the slowest wave the step admits, so the second axiom says that the longest wavelength available to a free particle is the circumference of the cosmological horizon. With the placement the momentum quantum $h/L\ell_P$ is $\hbar/R_\Lambda = 6.4\times10^{-61}$~kg\,m\,s$^{-1}$ and the energy quantum $\varepsilon_0 = \hbar c/R_\Lambda = 1.2\times10^{-33}$~eV, which is $h$ divided by the lap time $Lt_P = 2\pi R_\Lambda/c = 3.5\times10^{18}$~s and $2\pi k_B$ times the Gibbons--Hawking temperature~\cite{gibbons1977}.

The two determinations of $L$ are the two scales on which this circuit appears: from the horizon, $2\pi R_\Lambda/\ell_P = 6.4\times10^{61}$; from galaxies, $c^2/a_0\ell_P = 4.6\times10^{61}$ using $a_0 = c^2/L\ell_P$ of Ref.~\cite{korytko2026}. Only the product $L\ell_P$ enters $a_0$, so their agreement to 1.4 is Milgrom's $a_0 \approx c^2/2\pi R_\Lambda$ restated; a derivation of $a_0$ from the circuit alone would make the placement a prediction, and we do not have one.

The horizon has $A/\ell_P^2 = L^2/\pi$ pixels, a register being the $L$ along one great circle, and its Bekenstein--Hawking entropy~\cite{bekenstein1973} is
\begin{equation}
\frac{S}{k_B} = \frac{A}{4\ell_P^2} = \frac{L^2}{4\pi} = 3.3\times10^{122},
\label{eq:S}
\end{equation}
the value of Bousso's $N$-bound~\cite{bousso2000,bousso1999b}, which is the empty de Sitter horizon entropy written as $\pi R_\Lambda^2/\ell_P^2$. Equation~(\ref{eq:S}) is that number with $R_\Lambda = L\ell_P/2\pi$, an identity and not a derivation; what the register adds is that the integer fixing the log-dimension of a static patch is the $L$ that governs the qubit, Banks' proposal that positive $\Lambda$ and a finite Hilbert space go together~\cite{banks2000} with its integer named.

\section{Choosing a Direction}
\label{sec:direction}

The step translates, so it has a direction, but nothing so far has required one to be named: the algebra of Sec.~\ref{sec:algebra}, the bound of Sec.~\ref{sec:capacity}, and the horizon quantities of Sec.~\ref{sec:horizon} are the same along any direction. Dynamics is where a direction must be fixed, because a direction and its reverse make a line, and it is on a line that light has two components and matter has a chirality. Fix one. Section~\ref{sec:horizon} placed the ring on a great circle of the horizon, and every great circle has the same circumference, so which one the direction selects changes nothing and $L$ is unaffected. Along the direction the same $\mathbb{Z}_L$ carries the momenta $\kappa$ of Eq.~(\ref{eq:pk}), and the basis $|j\rangle$ conjugate to them is what position in whole cells means below: a coordinate on the register, periodic with the circumference because no wave is longer, and not a pixelation of the line. What follows is what the postulate gives on one line; several particles in several planes need a three-dimensional space, which is not constructed here and whose dimension Ref.~\cite{korytko2026d} derives.

\subsection{Light}
\label{sec:light}

For a massless mode $E = |p|c$, so by Eq.~(\ref{eq:pk})
\begin{equation}
E_\kappa = |\kappa|\,\varepsilon_0, \qquad \varepsilon_0 = \frac{hc}{L\ell_P} = \frac{2\pi\hbar}{L\,t_P},
\label{eq:eps0}
\end{equation}
and the phase advance per tick, $E_\kappa t_P/\hbar = 2\pi|\kappa|/L$, equals in magnitude the phase change across one cell, so the wave moves one cell per tick in the direction of its momentum. The wave $e^{i(kx-\omega t)}$ is invariant under one cell and one tick together. In operator form, $U(t_P) = e^{-iHt_P/\hbar}$ with $H = |\hat p|c$ acts as $X$ on every right-moving mode and $X^{-1}$ on every left-moving one: each tick, each component of light advances one cell in its own direction, and a mode with $\kappa$ wavelengths around the ring gains $\kappa$ phase steps in doing so.

A massless mode is a permitted state at every tick: its clock and its ruler are the same step, which is the postulate at its sharpest. A superposition of massless modes, a standing wave of light for instance, is also permitted at every tick, because every mode's phase advances by a permitted amount per tick. A massless mode never reverses direction; the term that reverses direction is the mass term of Sec.~\ref{sec:matter}, and on the register it is the only difference between light and matter.

\subsection{Matter}
\label{sec:matter}

Let $|{\rightarrow}\rangle$ and $|{\leftarrow}\rangle$ be the right- and left-moving light-like components, the two chiralities of a particle in one dimension, and write the state (\ref{eq:qubit}) in that basis. The velocity operator is $c\sigma_z$, and its expectation is
\begin{equation}
\langle v\rangle = c\,(P_\rightarrow - P_\leftarrow) = c\Big(\frac{2m}{L} - 1\Big) = c\cos\theta .
\label{eq:velocity}
\end{equation}
Palmer's count sets the velocity: in the string, the $m$ bits marked $+1$ are the right-movers and the $L-m$ marked $-1$ are the left-movers, the count is the probability of moving right, and the velocity is the excess of right-movers over left-movers in units of $c/L$. This right--left doublet is one qubit per particle, its chirality, the factor of two in a Dirac particle's Hilbert space that multiplies the $L$-dimensional orbital space of Sec.~\ref{sec:capacity}; the components at each momentum are that one qubit resolved momentum by momentum. The poles are light moving right and left. The equator, with $|{\rightarrow}\rangle$ and $|{\leftarrow}\rangle$ carrying momenta $\pm\kappa$, is a standing wave at rest, the interference pattern of the two components, with maxima every half wavelength; the shift $\phi$ is where those maxima sit, and one shift moves the pattern by one part in $L$ of its period.

Mass $M$ couples the two chiralities,
\begin{equation}
H = c\,\hat p\,\sigma_z + Mc^2\,\sigma_x, \qquad E = \pm\sqrt{p^2c^2 + M^2c^4},
\label{eq:dirac}
\end{equation}
the Dirac equation in one dimension, in which the mass term mixes the two chiralities at each momentum and the eigenstate has $\cos\theta = pc/E = v/c$. The $\sigma_x$ term reverses a quantum's direction, with angle per tick
\begin{equation}
\frac{Mc^2\,t_P}{\hbar} = \frac{M}{m_P} ,
\label{eq:flip}
\end{equation}
$4\times10^{-23}$ for an electron, one radian per Compton time $\hbar/Mc^2$. A massive particle at rest is light of Compton wavelength running right and left in equal measure, a right-mover becoming a left-mover in a quarter Compton period and back in the next; its rest energy is that of the trapped light. The lattice form of this construction, a particle moving at $c$ on a spacetime lattice and reversing with amplitude proportional to $M$ per step, is Feynman's checkerboard~\cite{feynman1965,jacobson1984,bialynicki1994}, whose continuum limit is the Dirac equation; the register is the checkerboard with the spacing set to $\ell_P$. In operator form the tick with mass is
\begin{equation}
U(t_P) = e^{-i(M/m_P)\sigma_x}\,X^{\sigma_z},
\label{eq:tick}
\end{equation}
the massless tick of Sec.~\ref{sec:light} followed by the reversal of Eq.~(\ref{eq:flip}). For $p\ell_P \ll \hbar$ and $M \ll m_P$ it is $1 - (it_P/\hbar)(c\hat p\sigma_z + Mc^2\sigma_x)$, which is Eq.~(\ref{eq:dirac}); its exact eigenphases obey
\begin{equation}
\cos\!\Big(\frac{E\,t_P}{\hbar}\Big) = \cos\!\Big(\frac{Mc^2\,t_P}{\hbar}\Big)\cos\!\Big(\frac{p\,\ell_P}{\hbar}\Big),
\label{eq:walk}
\end{equation}
which is $E = |p|c$ for $M = 0$ at every momentum, and Sec.~\ref{sec:pred} gives its expansion for matter.

Two consequences follow. First, the right--left reading of matter is a permitted setting only at intervals, the state itself being permitted at every tick. Equation~(\ref{eq:raqm}) is read in two ways, Sec.~\ref{sec:grid}: a mode of definite momentum is a record at every tick, its pattern that of Eq.~(\ref{eq:pk}) and its only advancing phase global and hidden, and what the grid constrains is the setting in which its right--left qubit may be read. Under Eq.~(\ref{eq:dirac}) the right--left qubit of a state of definite momentum turns about the axis $\hat n = (Mc^2, 0, pc)/E$ of the Bloch sphere, the $x$ axis at rest, at the rate $2E/\hbar$. The grid of Eq.~(\ref{eq:raqm}), rational $\cos\theta$ and rational $\phi/2\pi$, is invariant under rotations about $z$ by permitted angles and under rotations by $\pi$ about equatorial axes at permitted azimuths, which send $(\theta,\phi)$ to $(\pi-\theta, 2\alpha-\phi)$, and not under a rotation about $\hat n$ by a general angle. One tick of the mass term therefore takes a right--left superposition written on the grid off it, and the right--left reading is not a permitted setting until it returns. At rest $\hat n$ is equatorial, and the reading is permitted again once the Bloch vector has turned by $\pi$, after $\pi\hbar/2Mc^2$, a quarter Compton period, which is the time for a right-mover to become a left-mover. The register catches the direction of matter the way a stroboscope catches a wheel, only at the instants a spoke is at the mark.

In motion $\hat n$ tilts out of the equator, the half-turn lands off the grid, and the return waits for the full turn, $\pi\hbar/E$; the eigenstate itself, with $\cos\theta = pc/E$, is never on the right--left grid, so the velocity $\pm c$ of a moving massive particle is not a permitted reading at any tick, while its mean velocity $pc^2/E$, an orbital quantity, is. A superposition of momenta returns to the grid as a whole only if its spectrum is commensurate, which the register's own dispersion, Eq.~(\ref{eq:walk}), is not, apart from exceptional momenta: a wave packet of matter never returns to the grid as a whole, each component returning at its own time. In the nonrelativistic limit, $\hat p^2/2M$ on the ring, the spectrum is commensurate and every phase is again a permitted step after $ML\ell_P^2/\pi\hbar$, $4.6\times10^{-5}$~s for an electron; a nearest-neighbor hopping dispersion never returns at all, by Niven's theorem.

None of this touches $L$: the momentum quantum, the zone edge, and $\varepsilon_0$ are kinematic statements about the ring; the capacity follows from Born's rule; all are the same for every dispersion; and $L$ itself is frame-independent because every observer's horizon has the same radius. What the register does not settle is the Lorentz invariance of the lattice: a lattice of cells picks out a frame, the checkerboard recovers Lorentz invariance only in its continuum limit, and Palmer defers the question to future work~\cite{palmer2026c}.

Second, the global phase of a massive mode at rest runs at $Mc^2/\hbar$, which is $McL\ell_P/2\pi\hbar\approx4\times10^{38}$ steps per tick for an electron. That rate is not itself the obstacle; a photon of Compton momentum has the same phase rate and moves one cell per tick. The obstacle is that a global phase has no gradient: it advances by the same amount at every cell, so it displaces nothing. In Palmer's construction the global phase is $\xi$, a permutation with its origin fixed at creation and hidden~\cite{palmer2026c}. The register says why it must be: $\xi$ is the part of the rotation that is not motion, the particle's internal clock, and a uniform phase is invisible to every relational observable.


\subsection{Beyond one line}
\label{sec:beyond}

The two-mode results hold in any number of dimensions, since a superposition of $\mathbf{k}_1$ and $\mathbf{k}_2$ depends on position only through $\Delta\mathbf{k}\cdot\mathbf{x}$, and the translation reading holds as written for collinear modes. Beyond that the free dynamics is planar and exact in the mode basis. A planar mode, labeled by its magnitude $\kappa$ and its direction, advances $|\kappa|$ phase steps per tick whatever its direction, so a wave moves one Planck length along its own $\hat{\mathbf{k}}$ per tick. A massive mode along $\hat{\mathbf{k}}$ carries its right--left doublet along $\hat{\mathbf{k}}$ and obeys Sec.~\ref{sec:matter} in that direction.

The two-mode dictionary of Sec.~\ref{sec:matter} holds for any pair of modes in any dimension: the count is the probability of one mode, and one shift moves their interference pattern by one part in $L$ of its period along $\hat{\Delta\mathbf{k}}$. What is one-dimensional is only the whole-cell version of that reading: one cell of translation is a whole number of shifts when $|\Delta\mathbf{k}|$ is a multiple of $2\pi/L\ell_P$, which it always is along the ring and, Pythagorean pairs aside, is not for non-collinear modes.

Every elementary process is planar: two colliding momenta span a plane; a vertex's three momenta, three legs being what Ref.~\cite{korytko2026d} derives, sum to zero and lie in one; and a boost within the plane keeps them there. A planar momentum is a ring mode for its magnitude and a point on one great circle for its direction, $L^2$ labels, and the horizon has $L^2/\pi$ pixels; to within $\pi$ the horizon as it stands labels the momenta of every elementary process. A state of several particles in several planes needs of order $L^3$ labels, and the degrees of freedom the horizon affords it number $L^2/4\pi$, Eq.~(\ref{eq:S}). That number does not compete with Eq.~(\ref{eq:nmaxspan}) and is not commensurable with it: the first counts degrees of freedom that exist, the second the outcomes to which one pure state can give cells, and nothing requires a state to occupy every degree of freedom that exists. How such a bulk state is written on the horizon is the open problem of de Sitter holography~\cite{bousso1999,banksfischler2001,parikh2003}, which this paper does not address, and interactions are constructed in no dimension here. Ref.~\cite{korytko2026d} adds a clause on the phase the step carries across each link between cells; the interaction it admits emits or absorbs a massless quantum, and decays of matter are left to a later paper. Gravity is untouched: the Jacobson and Verlinde routes of Ref.~\cite{korytko2026} need only the pixel on every screen, never a state.

\section{Predictions and Falsifiers}
\label{sec:pred}

The register inherits the predictions of Ref.~\cite{korytko2026}, corrects one criterion, and adds no test.

\textit{Inherited.} The two predictions of Ref.~\cite{korytko2026} stand, since this paper rests on the same $L$. First, a state written on $N$ records gives cells to at most $L$ joint outcomes, Eq.~(\ref{eq:nmaxspan}): a ceiling of about 200 usefully entangled qubits, which Sec.~\ref{sec:capacity} places between 203 and 205 and does not claim is attained, the $\log_2 L$ of Ref.~\cite{korytko2026} in place of the 212 it called exact. The ceiling holds for any physical realization, $M$ bosonic modes with photon-number cutoff $n$ reaching it at $n^M \approx L$. Operationally: run a random circuit on $N$ fault-tolerant logical qubits, deep enough to anticoncentrate, which is logarithmic depth~\cite{dalzell2022}, and then its inverse. In standard quantum mechanics the state returns to the initial one with the gate fidelity. On the register the return probability falls to about $(L/2^N)^2$ for $N$ beyond $\lfloor\log_2 L\rfloor$, since the intermediate state has $2^N$ branches and the register has cells for $L$. The same circuit built from Clifford gates, whose intermediate states have at most three blocks to a record, returns at any $N$. The test needs no classical simulation, only the check of a return; no such machine exists. Second, $\Lambda$ is constant, so the dark-energy equation of state is $w = -1$ exactly, against the DESI DR2 preference for evolving dark energy~\cite{desi2025}, which the framework predicts will not survive.

\textit{Definite, but not discriminating.} Momentum is quantized in units of $\hbar/R_\Lambda$ and bounded at $\pi\hbar/\ell_P$. Light is a record at every Planck tick, and the right--left reading of matter at rest is a permitted setting once a quarter Compton period. The settings in which any two-mode subspace of a free particle may be read are restricted by Palmer's Impossible Triangle Corollary. The maximum group velocity of a particle of mass $M$ is $c\cos(M/m_P)$, reached at $p = \pi\hbar/2\ell_P$ and below $c$ by $(M/m_P)^2/2 = 9\times10^{-46}$ for an electron, the group velocity of matter vanishing at the zone edge while light's is $c$ at every momentum. The lattice fixes what Lorentz violation the register entails: the tick's exact dispersion, Eq.~(\ref{eq:walk}), is $E = |p|c$ for light at every momentum and, for matter,
\begin{equation}
E^2 = p^2c^2 + M^2c^4 - \tfrac{1}{3}\,M^2c^4\Big(\frac{p\,\ell_P}{\hbar}\Big)^2 + \dots ,
\label{eq:lattice}
\end{equation}
a fractional correction, of order $(Mc^2/E_P)^2 \approx 6\times10^{-39}$ for a proton, that does not grow with energy. The register therefore predicts no energy-dependent speed of light at any order and would be falsified by one; the correction for matter is beyond any measurement. The lattice likewise deforms the commutator: with the lattice momentum $\hat p_{\rm lat} = (\hbar/\ell_P)\sin(\hat p\ell_P/\hbar)$, $[\hat x,\hat p_{\rm lat}] = i\hbar\cos(\hat p\ell_P/\hbar) \approx i\hbar\,[1 - \tfrac12(\hat p\ell_P/\hbar)^2]$, a generalized uncertainty principle of the maximal-momentum type with coefficient $\beta_0 = -1/2$ for one particle. Each constituent of a body carries its own record, so for the center of mass of $N$ of them the coefficient is $-1/2N^2$, the suppression Amelino-Camelia argued any composite must show~\cite{amelino2013}. The mechanical-resonator bounds, $|\beta_0| < 3\times10^7$~\cite{bawaj2015}, $5\times10^6$~\cite{bushev2019}, and now 62 from a cryogenic quartz resonator~\cite{campbell2023}, by the route Pikovski et al. proposed~\cite{pikovski2012}, are on center-of-mass modes of $N \approx 10^{20}$ constituents: two orders of magnitude above one particle's coefficient and forty above the body's. Because a two-path superposition of a composite's center of mass is one qubit whatever the mass, and the register has no collapse mechanism, it predicts no intrinsic, mass-dependent loss of matter-wave contrast, in agreement with experiment to $2.5\times10^4$ atomic mass units~\cite{fein2019} and unlike collapse models. None of these can distinguish the register from standard quantum mechanics.

The register's intrinsic scales are $m_P(2\pi/L)^{1/n}$: $1.2\times10^{-33}$ eV, $3.8$ meV, $56$ MeV, and $6.8$ TeV. The second is, to a factor 1.7, the dark-energy scale $\rho_\Lambda^{1/4}$, which follows from Eq.~(\ref{eq:closure}) as $(3\pi/2)^{1/4}E_P/\sqrt L = 2.2$ meV, the fourth root of the $L^2$ suppression of the Planck density counted in Ref.~\cite{korytko2026b}; the third is Weinberg's $(\hbar^2H_0/Gc)^{1/3}$~\cite{weinberg1972}. That the neutrino and the pion sit near the second and third rungs is the set of large-number coincidences known since Dirac, in one integer; the register offers no mechanism for them.

\section{What the Placement Costs}
\label{sec:cost}

Placing every register on the horizon can mean two things. The first, that the cells have a definite orientation in space, costs nothing. A cyclic register has no pole and no first cell; both are coordinates on the Bloch sphere, and every quantity in Eq.~(\ref{eq:raqm}) is relational, an overlap of two states or a phase relative to a reference, which Sec.~\ref{sec:matter} shows to be a position relative to a reference. Palmer's Bell conditions are on angles between settings~\cite{palmer2026b}, which are rotation-invariant; he fixes the gauge by placing the pole at his middle setting~\cite{palmer2026c}, and his distinction between a nominal setting and the exact one within it~\cite{palmer2024} says the gauge cannot be probed.

The postulate's last clause carries this from the step to the states written on it. Not only is there no first cell; there is no phase reference outside a state, its global phase being compared with nothing. Nothing in Eq.~(\ref{eq:raqm}) reads a state's global phase and no dynamics above does either, so the extension costs nothing and closes the one place where an unprobed convention could have passed for data. It has no effect on the capacity: a state's global phase is the origin of its cell labels, which no count and no block reads.

The second, that different observers' horizons are the same cells, is de Sitter complementarity~\cite{susskind1993,bousso1999,dyson2002,banksfischler2001,parikh2003}: distinct static patches describe one finite set of degrees of freedom, not independent copies. The shared reading follows from unitarity within a patch by the no-cloning argument used for black holes~\cite{susskind1994,wootters1982}, taken over without proof; whether that unitarity is consistent with a smooth horizon is the firewall question~\cite{amps2013}, open in de Sitter and not addressed here. Nothing in Secs.~\ref{sec:register} to \ref{sec:pred} depends on it. A finite Hilbert space of log-dimension $L^2/4\pi$ entails recurrences and a low-entropy past that requires explanation~\cite{dyson2002}, to which Palmer's nonergodic Invariant Set Postulate~\cite{palmer2009} is RaQM's response; it entails positive $\Lambda$~\cite{banks2000} and no exact asymptotic observables~\cite{witten2001}; and finite-dimensional Hilbert spaces for quantum gravity have been advocated independently~\cite{buniy2005,bao2017,carroll2023}.

A finite register also says what it does not hold. The bound of Sec.~\ref{sec:capacity} is on a pure state written on one register. An observer outside a black hole describes it by a mixed state, whose entropy counts the records the description ranges over, as Eq.~(\ref{eq:S}) counts them for the cosmological horizon; Sec.~\ref{sec:capacity} counts the outcomes to which one record can give cells, and the two are not the same quantity. What the register does not admit is a scrambled pure state above $\lfloor\log_2 L\rfloor$ qubits, and the pure state of a hole together with its radiation, which unitarity of the pair asserts and which scrambling makes generic~\cite{sekino2008}, is one. No observer's register writes it. That is the register's position on the information question, recorded here and not resolved; it agrees with there being no exact observables in de Sitter space~\cite{witten2001}.

\section{Discussion}
\label{sec:discussion}

The rational states Ref.~\cite{korytko2026} took from Palmer and the universality it had to stipulate are one statement about one ring. A string of $L$ cells turns at the rate its energy sets. Counted, the pattern is a Born probability; for the right- and left-moving components that probability sets the velocity. Shifted, it is a phase; on the ring read as space, one cell is a translation by one Planck length, the slowest wave gains one step of phase, Palmer's $2\pi/L$, per translation, and for light the phase and the translation are the same thing at every tick. Mass is the amplitude per tick to reverse direction, a particle at rest is a standing wave of trapped light, and the Compton clock is the global phase that Palmer hides and that, having no gradient, displaces nothing. The circumference is the longest wavelength a particle can have, and placing it on the horizon is the second axiom with that content attached.

Less is stipulated than in Ref.~\cite{korytko2026}, though not less is assumed. Its first axiom is a consequence of the postulate, Eq.~(\ref{eq:raqm}) with $L$ universal, since the cell belongs to space rather than to any system in it; its second, the placement of the ring on the horizon, is taken over here and derived in Ref.~\cite{korytko2026d}, and it is what gives $L$ its value, Eq.~(\ref{eq:ident}). What changes is what must be defended: a universal $L$ no longer has to be imposed against Palmer's mass-dependent one, because a grain that varies with the system cannot be the grain that does not. The other gains are structural. RaQM's phase step is momentum quantization on a closed ring, its qubit is the right--left doublet of a Dirac particle, and its hidden global phase is the Compton clock. Its capacity bound, two cells to every branch of a state, is $\lfloor\log_2 L\rfloor$, soft by a unit in resolution: a ceiling of about 200 that no representation may exceed and that holds whatever the qubits are made of. What remains open is what was open before, the coefficient in $a_0$ and complementarity; the three-dimensional bulk is the subject of Ref.~\cite{korytko2026d}. Ref.~\cite{korytko2026} remains the statement the framework rests on, its exact capacity criterion corrected here. Palmer has left the discretization of space and time to a future paper and announced work on vacuum energy~\cite{palmer2026c}; this paper and Ref.~\cite{korytko2026b} attempt those questions with a universal $L$.


\begin{acknowledgments}
The author used an AI assistant (Anthropic's Claude) for literature checking, numerical verification, and editing of the manuscript. The physical proposal, the postulate, and the claims are the author's own.
\end{acknowledgments}

\begin{thebibliography}{99}
\bibitem{palmer2020} T. N. Palmer, Discretization of the Bloch sphere, fractal invariant sets and Bell's theorem, Proc. R. Soc. A \textbf{476}, 20190350 (2020). \url{https://doi.org/10.1098/rspa.2019.0350}
\bibitem{palmer2026} T. N. Palmer, Rational quantum mechanics: Testing quantum theory with quantum computers, Proc. Natl. Acad. Sci. USA \textbf{123}, e2523350123 (2026). \url{https://doi.org/10.1073/pnas.2523350123}; arXiv:2510.02877.
\bibitem{palmer2026b} T. N. Palmer, Impossible counterfactuals, discrete Hilbert space and Bell's theorem, J. Phys.: Conf. Ser. \textbf{3189}, 012006 (2026). \url{https://doi.org/10.1088/1742-6596/3189/1/012006}; arXiv:2601.14941.
\bibitem{palmer2024} T. N. Palmer, Superdeterminism without conspiracy, Universe \textbf{10}, 47 (2024). \url{https://doi.org/10.3390/universe10010047}
\bibitem{palmer2026c} T. N. Palmer, Solving the mysteries of quantum mechanics: why nature abhors a continuum, arXiv:2602.16382v1 (2026).
\bibitem{korytko2026} A. Korytko, Gravitation from Hilbert-space granularity: pinning the parameter $L$ of rational quantum mechanics across scales (2026). \url{https://doi.org/10.5281/zenodo.22255462}
\bibitem{weyl1931} H. Weyl, \emph{The Theory of Groups and Quantum Mechanics} (Methuen, London, 1931; Dover reprint, New York, 1950), Ch.~IV, Sec.~14.
\bibitem{schwinger1960} J. Schwinger, Unitary operator bases, Proc. Natl. Acad. Sci. USA \textbf{46}, 570 (1960). \url{https://doi.org/10.1073/pnas.46.4.570}
\bibitem{vourdas2004} A. Vourdas, Quantum systems with finite Hilbert space, Rep. Prog. Phys. \textbf{67}, 267 (2004). \url{https://doi.org/10.1088/0034-4885/67/3/R03}
\bibitem{massar2008} S. Massar and P. Spindel, Uncertainty relation for the discrete Fourier transform, Phys. Rev. Lett. \textbf{100}, 190401 (2008). \url{https://doi.org/10.1103/PhysRevLett.100.190401}
\bibitem{donoho1989} D. L. Donoho and P. B. Stark, Uncertainty principles and signal recovery, SIAM J. Appl. Math. \textbf{49}, 906 (1989). \url{https://doi.org/10.1137/0149053}
\bibitem{feynman1965} R. P. Feynman and A. R. Hibbs, \emph{Quantum Mechanics and Path Integrals} (McGraw-Hill, New York, 1965), pp.~34--36, Problem 2-6.
\bibitem{jacobson1984} T. Jacobson and L. S. Schulman, Quantum stochastics: the passage from a relativistic to a non-relativistic path integral, J. Phys. A: Math. Gen. \textbf{17}, 375 (1984). \url{https://doi.org/10.1088/0305-4470/17/2/023}
\bibitem{bialynicki1994} I. Bialynicki-Birula, Weyl, Dirac, and Maxwell equations on a lattice as unitary cellular automata, Phys. Rev. D \textbf{49}, 6920 (1994). \url{https://doi.org/10.1103/PhysRevD.49.6920}
\bibitem{weinberg1972} S. Weinberg, \emph{Gravitation and Cosmology} (Wiley, New York, 1972), pp.~619--620.
\bibitem{korytko2026b} A. Korytko, The vacuum catastrophe and the granularity of Hilbert space (2026). \url{https://doi.org/10.5281/zenodo.22683543}
\bibitem{korytko2026d} A. Korytko, The dimension of space from the granularity of Hilbert space (2026). \url{https://doi.org/10.5281/zenodo.22678537}
\bibitem{pikovski2012} I. Pikovski, M. R. Vanner, M. Aspelmeyer, M. S. Kim, and \v{C}. Brukner, Probing Planck-scale physics with quantum optics, Nat. Phys. \textbf{8}, 393 (2012). \url{https://doi.org/10.1038/nphys2262}
\bibitem{amelino2013} G. Amelino-Camelia, Challenge to macroscopic probes of quantum spacetime based on noncommutative geometry, Phys. Rev. Lett. \textbf{111}, 101301 (2013). \url{https://doi.org/10.1103/PhysRevLett.111.101301}
\bibitem{bawaj2015} M. Bawaj, C. Biancofiore, M. Bonaldi, F. Bonfigli, A. Borrielli, G. Di Giuseppe, L. Marconi, F. Marino, R. Natali, A. Pontin, G. A. Prodi, E. Serra, D. Vitali, and F. Marin, Probing deformed commutators with macroscopic harmonic oscillators, Nat. Commun. \textbf{6}, 7503 (2015). \url{https://doi.org/10.1038/ncomms8503}
\bibitem{bushev2019} P. A. Bushev, J. Bourhill, M. Goryachev, N. Kukharchyk, E. Ivanov, S. Galliou, M. E. Tobar, and S. Danilishin, Testing the generalized uncertainty principle with macroscopic mechanical oscillators and pendulums, Phys. Rev. D \textbf{100}, 066020 (2019). \url{https://doi.org/10.1103/PhysRevD.100.066020}
\bibitem{campbell2023} W. M. Campbell, M. E. Tobar, M. Goryachev, and S. Galliou, Improved constraints on minimum length models with a macroscopic low loss phonon cavity, Phys. Rev. D \textbf{108}, 102006 (2023). \url{https://doi.org/10.1103/PhysRevD.108.102006}
\bibitem{fein2019} Y. Y. Fein, P. Geyer, P. Zwick, F. Kia\l{}ka, S. Pedalino, M. Mayor, S. Gerlich, and M. Arndt, Quantum superposition of molecules beyond 25 kDa, Nat. Phys. \textbf{15}, 1242 (2019). \url{https://doi.org/10.1038/s41567-019-0663-9}
\bibitem{sekino2008} Y. Sekino and L. Susskind, Fast scramblers, J. High Energy Phys. \textbf{10}, 065 (2008). \url{https://doi.org/10.1088/1126-6708/2008/10/065}; arXiv:0808.2096.
\bibitem{dalzell2022} A. M. Dalzell, N. Hunter-Jones, and F. G. S. L. Brand\~ao, Random quantum circuits anticoncentrate in log depth, PRX Quantum \textbf{3}, 010333 (2022). \url{https://doi.org/10.1103/PRXQuantum.3.010333}; arXiv:2011.12277.
\bibitem{niven1956} I. Niven, \emph{Irrational Numbers} (Mathematical Association of America, Washington, DC, 1956).
\bibitem{jahnel2010} J. Jahnel, When is the (co)sine of a rational angle equal to a rational number?, arXiv:1006.2938 (2010).
\bibitem{bousso2000} R. Bousso, Positive vacuum energy and the $N$-bound, J. High Energy Phys. \textbf{11}, 038 (2000). \url{https://doi.org/10.1088/1126-6708/2000/11/038}; hep-th/0010252.
\bibitem{bousso1999} R. Bousso, Quantum global structure of de Sitter space, Phys. Rev. D \textbf{60}, 063503 (1999). \url{https://doi.org/10.1103/PhysRevD.60.063503}
\bibitem{bousso1999b} R. Bousso, A covariant entropy conjecture, J. High Energy Phys. \textbf{07}, 004 (1999). \url{https://doi.org/10.1088/1126-6708/1999/07/004}
\bibitem{jacobson1995} T. Jacobson, Thermodynamics of spacetime: the Einstein equation of state, Phys. Rev. Lett. \textbf{75}, 1260 (1995). \url{https://doi.org/10.1103/PhysRevLett.75.1260}
\bibitem{susskind1993} L. Susskind, L. Thorlacius, and J. Uglum, The stretched horizon and black hole complementarity, Phys. Rev. D \textbf{48}, 3743 (1993). \url{https://doi.org/10.1103/PhysRevD.48.3743}
\bibitem{susskind1994} L. Susskind and L. Thorlacius, Gedanken experiments involving black holes, Phys. Rev. D \textbf{49}, 966 (1994). \url{https://doi.org/10.1103/PhysRevD.49.966}
\bibitem{banks2000} T. Banks, Cosmological breaking of supersymmetry, or little Lambda goes back to the future II, arXiv:hep-th/0007146 (2000).
\bibitem{witten2001} E. Witten, Quantum gravity in de Sitter space, arXiv:hep-th/0106109 (2001).
\bibitem{buniy2005} R. V. Buniy, S. D. H. Hsu, and A. Zee, Is Hilbert space discrete?, Phys. Lett. B \textbf{630}, 68 (2005). \url{https://doi.org/10.1016/j.physletb.2005.09.084}; arXiv:hep-th/0508039.
\bibitem{bao2017} N. Bao, S. M. Carroll, and A. Singh, The Hilbert space of quantum gravity is locally finite-dimensional, Int. J. Mod. Phys. D \textbf{26}, 1743013 (2017). \url{https://doi.org/10.1142/S0218271817430131}; arXiv:1704.00066.
\bibitem{carroll2023} S. M. Carroll, Completely discretized, finite quantum mechanics, Found. Phys. \textbf{53}, 90 (2023). \url{https://doi.org/10.1007/s10701-023-00726-6}
\bibitem{palmer2009} T. N. Palmer, The Invariant Set Postulate: a new geometric framework for the foundations of quantum theory and the role played by gravity, Proc. R. Soc. A \textbf{465}, 3165 (2009). \url{https://doi.org/10.1098/rspa.2009.0080}
\bibitem{desi2025} DESI Collaboration, M. Abdul-Karim \emph{et al.}, DESI DR2 results. II. Measurements of baryon acoustic oscillations and cosmological constraints, Phys. Rev. D \textbf{112}, 083515 (2025). \url{https://doi.org/10.1103/tr6y-kpc6}; arXiv:2503.14738.
\bibitem{amps2013} A. Almheiri, D. Marolf, J. Polchinski, and J. Sully, Black holes: complementarity or firewalls?, J. High Energy Phys. \textbf{02}, 062 (2013). \url{https://doi.org/10.1007/JHEP02(2013)062}; arXiv:1207.3123.
\bibitem{wootters1982} W. K. Wootters and W. H. Zurek, A single quantum cannot be cloned, Nature \textbf{299}, 802 (1982). \url{https://doi.org/10.1038/299802a0}
\bibitem{dyson2002} L. Dyson, M. Kleban, and L. Susskind, Disturbing implications of a cosmological constant, J. High Energy Phys. \textbf{10}, 011 (2002). \url{https://doi.org/10.1088/1126-6708/2002/10/011}; hep-th/0208013.
\bibitem{banksfischler2001} T. Banks and W. Fischler, M-theory observables for cosmological space-times, arXiv:hep-th/0102077 (2001).
\bibitem{parikh2003} M. K. Parikh, I. Savonije, and E. Verlinde, Elliptic de Sitter space: dS/$\mathbb{Z}_2$, Phys. Rev. D \textbf{67}, 064005 (2003). \url{https://doi.org/10.1103/PhysRevD.67.064005}
\bibitem{thooft1993} G. 't Hooft, Dimensional reduction in quantum gravity, arXiv:gr-qc/9310026 (1993).
\bibitem{susskind1995} L. Susskind, The world as a hologram, J. Math. Phys. \textbf{36}, 6377 (1995). \url{https://doi.org/10.1063/1.531249}
\bibitem{bekenstein1973} J. D. Bekenstein, Black holes and entropy, Phys. Rev. D \textbf{7}, 2333 (1973). \url{https://doi.org/10.1103/PhysRevD.7.2333}
\bibitem{gibbons1977} G. W. Gibbons and S. W. Hawking, Cosmological event horizons, thermodynamics, and particle creation, Phys. Rev. D \textbf{15}, 2738 (1977). \url{https://doi.org/10.1103/PhysRevD.15.2738}
\bibitem{verlinde2011} E. Verlinde, On the origin of gravity and the laws of Newton, J. High Energy Phys. \textbf{04}, 029 (2011). \url{https://doi.org/10.1007/JHEP04(2011)029}
\bibitem{milgrom1983} M. Milgrom, A modification of the Newtonian dynamics as a possible alternative to the hidden mass hypothesis, Astrophys. J. \textbf{270}, 365 (1983). \url{https://doi.org/10.1086/161130}
\bibitem{mcgaugh2016} S. S. McGaugh, F. Lelli, and J. M. Schombert, Radial acceleration relation in rotationally supported galaxies, Phys. Rev. Lett. \textbf{117}, 201101 (2016). \url{https://doi.org/10.1103/PhysRevLett.117.201101}
\bibitem{planck2018} Planck Collaboration, Planck 2018 results. VI. Cosmological parameters, Astron. Astrophys. \textbf{641}, A6 (2020). \url{https://doi.org/10.1051/0004-6361/201833910}
\end{thebibliography}

\end{document}
